\documentclass[12pt,a4paper]{article}

\usepackage[margin=1in]{geometry}
\usepackage{amsmath, amssymb, amsthm}
\usepackage{enumitem}
\usepackage{booktabs}
\usepackage{tikz}
\usetikzlibrary{arrows.meta, positioning}

\setlength{\parindent}{0pt}
\setlength{\parskip}{6pt}

\newtheoremstyle{notestyle}
  {8pt}
  {8pt}
  {\normalfont}
  {}
  {\bfseries}
  {.}
  {0.5em}
  {}

\theoremstyle{notestyle}
\newtheorem{definition}{Definition}
\newtheorem{example}{Example}

\newenvironment{answer}
{\par\textbf{Answer.}\ }
{\hfill$\square$\par}

\newcommand{\Prob}{\mathbb{P}}
\newcommand{\E}{\mathbb{E}}

\title{\textbf{Concept Clarifications: Markov Chain Calculations}}
\author{MA4034: Time Series and Stochastic Processes}
\date{}

\begin{document}
\maketitle

\section{Why Does $\lambda^n$ Appear in $P^n$?}

Suppose $P$ is a transition matrix. The $n$-step transition matrix is
\[
P^n.
\]

Directly multiplying
\[
P\cdot P\cdot P\cdots P
\]
many times is difficult. Eigenvalues help because they turn repeated matrix multiplication into repeated scalar multiplication.

\subsection*{Key Idea}

If
\[
Pv=\lambda v,
\]
then applying $P$ repeatedly gives
\[
P^2v=P(Pv)=P(\lambda v)=\lambda Pv=\lambda^2v.
\]
Similarly,
\[
P^3v=\lambda^3v.
\]
In general,
\[
P^nv=\lambda^nv.
\]

That is why powers of eigenvalues, such as
\[
\lambda^n,
\]
appear when calculating $P^n$.

\subsection*{Diagonalization}

If a matrix can be diagonalized, then
\[
P=RDR^{-1},
\]
where $D$ is a diagonal matrix containing the eigenvalues.

Then
\[
P^n=(RDR^{-1})^n.
\]

When we multiply this out,
\[
P^n=RD^nR^{-1}.
\]

This happens because the middle terms cancel:
\[
RDR^{-1}RDR^{-1}=RD^2R^{-1}.
\]

For a diagonal matrix,
\[
D=
\begin{pmatrix}
\lambda_1 & 0\\
0 & \lambda_2
\end{pmatrix},
\]
we have
\[
D^n=
\begin{pmatrix}
\lambda_1^n & 0\\
0 & \lambda_2^n
\end{pmatrix}.
\]

So the powers $\lambda_1^n,\lambda_2^n$ come from raising the diagonal matrix to the $n$th power.

\section{ON/OFF Chain: Meaning of the Limit}

For the ON/OFF chain,
\[
P=
\begin{pmatrix}
1-p & p\\
q & 1-q
\end{pmatrix}.
\]

Here:
\begin{itemize}[leftmargin=1.2cm]
    \item $p$ is the probability of switching from OFF to ON,
    \item $q$ is the probability of switching from ON to OFF.
\end{itemize}

The eigenvalues are
\[
\lambda_1=1,
\qquad
\lambda_2=1-p-q.
\]

The $n$-step transition matrix can be written as
\[
P^n=
\frac{1}{p+q}
\begin{pmatrix}
q & p\\
q & p
\end{pmatrix}
+
\frac{(1-p-q)^n}{p+q}
\begin{pmatrix}
p & -p\\
-q & q
\end{pmatrix}.
\]

The important part is
\[
(1-p-q)^n.
\]

If
\[
|1-p-q|<1,
\]
then
\[
\lim_{n\to\infty}(1-p-q)^n=0.
\]

So the second matrix disappears in the long run:
\[
\lim_{n\to\infty}P^n=
\frac{1}{p+q}
\begin{pmatrix}
q & p\\
q & p
\end{pmatrix}.
\]

Therefore,
\[
\lim_{n\to\infty}P^n=
\begin{pmatrix}
\frac{q}{p+q} & \frac{p}{p+q}\\
\frac{q}{p+q} & \frac{p}{p+q}
\end{pmatrix}.
\]

Both rows are the same. That means that after a long time, the probability of being OFF or ON does not depend on where we started.

\[
\Prob(\text{OFF in long run})=\frac{q}{p+q},
\]
\[
\Prob(\text{ON in long run})=\frac{p}{p+q}.
\]

\subsection*{Why is $\Prob(\text{OFF})=\frac{q}{p+q}$?}

This may feel reversed at first. But remember:

\[
q=\Prob(\text{ON}\to \text{OFF}).
\]

If $q$ is large, the system leaves ON quickly and goes to OFF often. So the system spends more time OFF.

Similarly,

\[
p=\Prob(\text{OFF}\to \text{ON}).
\]

If $p$ is large, the system leaves OFF quickly and goes to ON often. So the system spends more time ON.

\section{Series Criterion: What Do $I_n$, $S$, and $\E_i[S]$ Mean?}

Let
\[
I_n=
\begin{cases}
1, & X_n=i,\\
0, & X_n\neq i.
\end{cases}
\]

This is an indicator variable.

It means:
\begin{itemize}[leftmargin=1.2cm]
    \item $I_n=1$ if the chain is in state $i$ at time $n$,
    \item $I_n=0$ if the chain is not in state $i$ at time $n$.
\end{itemize}

Then
\[
S=\sum_{n=1}^{\infty}I_n
\]
counts how many times the chain returns to state $i$ after time 0.

\subsection*{Why is $\E_i[I_n]=\Prob_i(X_n=i)$?}

For an indicator variable,
\[
\E[I_n]=1\cdot \Prob(I_n=1)+0\cdot \Prob(I_n=0).
\]

So
\[
\E[I_n]=\Prob(I_n=1).
\]

But
\[
I_n=1
\]
means
\[
X_n=i.
\]

Therefore,
\[
\E_i[I_n]=\Prob_i(X_n=i)=p_{ii}^{(n)}.
\]

Hence,
\[
\E_i[S]
=
\sum_{n=1}^{\infty}\E_i[I_n]
=
\sum_{n=1}^{\infty}p_{ii}^{(n)}.
\]

\subsection*{Meaning}

\[
\sum_{n=1}^{\infty}p_{ii}^{(n)}
\]
is the expected total number of returns to state $i$.

If this sum is finite, state $i$ is transient.

If this sum is infinite, state $i$ is recurrent.

\section{Ergodic Chain: Why Every Row Becomes $\pi$}

For an ergodic Markov chain,
\[
\lim_{n\to\infty}P^n
\]
has every row equal to the stationary distribution $\pi$.

This means:

\[
p_{ij}^{(n)}\to \pi_j.
\]

The starting state is $i$, and the final state is $j$.

In the long run, the probability of ending in state $j$ becomes $\pi_j$, no matter where we started.

So if
\[
\pi=(0.2,0.5,0.3),
\]
then
\[
\lim_{n\to\infty}P^n
=
\begin{pmatrix}
0.2 & 0.5 & 0.3\\
0.2 & 0.5 & 0.3\\
0.2 & 0.5 & 0.3
\end{pmatrix}.
\]

Each row is the same because the chain forgets its starting state.

\section{Why $P(S)=P(S,S)+P(S,R)$?}

Suppose the weather can be sunny $(S)$ or rainy $(R)$ on two consecutive days.

The event
\[
\text{first day is sunny}
\]
can happen in two ways:
\[
(S,S)
\]
or
\[
(S,R).
\]

These are disjoint events. Therefore,
\[
\Prob(S)=\Prob(S,S)+\Prob(S,R).
\]

For example, if
\[
\Prob(S,S)=0.7
\]
and
\[
\Prob(S,R)=0.1,
\]
then
\[
\Prob(S)=0.7+0.1=0.8.
\]

\section{Insurance Example: Why Part (b) and Part (c) Are Different}

The insurance example has three states:
\[
H=\text{high risk},\qquad I=\text{intermediate risk},\qquad L=\text{low risk}.
\]

The question asks:

\begin{itemize}[leftmargin=1.2cm]
    \item[(b)] If a customer starts as low risk, how many years can they expect to stay there?
    \item[(c)] How many years pass on average between two consecutive visits to the high-risk category?
\end{itemize}

These are different types of questions.

\subsection{Part (b): How Long Do They Stay in Low Risk?}

This asks about the first time leaving state $L$.

If the customer is in low risk:
\begin{itemize}[leftmargin=1.2cm]
    \item with probability $0.8$, they stay low risk,
    \item with probability $0.2$, they leave low risk.
\end{itemize}

So the number of years spent in low risk before leaving has a geometric distribution with success probability
\[
0.2.
\]

Here ``success'' means leaving low risk.

Therefore, the expected number of years spent in low risk is
\[
\frac{1}{0.2}=5.
\]

\[
\boxed{\text{Expected time staying low risk}=5\text{ years}.}
\]

\subsection{Why Not Use $1/\pi_L$ for Part (b)?}

The value
\[
\frac1{\pi_L}
\]
means the average time between two visits to low risk in the long run.

But part (b) does not ask:

\[
\text{How long between visits to low risk?}
\]

It asks:

\[
\text{Starting in low risk, how long until we leave low risk?}
\]

That is a first exit time, not a recurrence time.

So we use
\[
\frac1{\Prob(\text{leave }L)}
=
\frac1{0.2}
=5.
\]

\subsection{Part (c): Time Between Consecutive Visits to High Risk}

Part (c) asks about recurrence time.

It asks:

\[
\text{High risk}\to \text{leave/move around}\to \text{High risk again}.
\]

For an ergodic Markov chain, the mean recurrence time of state $i$ is
\[
\E_i[\tau_i]=\frac1{\pi_i}.
\]

So for high risk,
\[
\E_H[\tau_H]=\frac1{\pi_H}.
\]

For the insurance example, the stationary distribution is
\[
\pi=\left(\frac5{29},\frac6{29},\frac{18}{29}\right)
\]
using state order
\[
(H,I,L).
\]

Therefore,
\[
\pi_H=\frac5{29}.
\]

So
\[
\E_H[\tau_H]
=
\frac1{\pi_H}
=
\frac1{5/29}
=
\frac{29}{5}
=5.8.
\]

\[
\boxed{\text{Average time between visits to high risk}=5.8\text{ years}.}
\]

\subsection{Main Difference}

\[
\begin{array}{c|c|c}
\text{Question} & \text{Meaning} & \text{Method}\\
\hline
\text{How long stay in }L? & \text{time until leaving }L & 1/\Prob(\text{leave }L)\\
\text{Time between visits to }H? & \text{return time to }H & 1/\pi_H
\end{array}
\]

\section{Practice Questions}

\begin{enumerate}[label=\textbf{Q\arabic*.}, leftmargin=1.2cm]
    \item For the ON/OFF chain, explain why the term $(1-p-q)^n$ disappears in the limit when $|1-p-q|<1$.

    \item Suppose
    \[
    \pi=(0.3,0.7).
    \]
    What will $\lim_{n\to\infty}P^n$ look like if the Markov chain is ergodic?

    \item If
    \[
    \Prob(S,S)=0.6,\qquad \Prob(S,R)=0.2,
    \]
    find $\Prob(S)$.

    \item A Markov chain has state $A$. From $A$, it stays in $A$ with probability $0.75$ and leaves $A$ with probability $0.25$. If the chain starts at $A$, what is the expected time spent in $A$ before leaving?

    \item In an ergodic Markov chain, suppose the stationary probability of state $B$ is
    \[
    \pi_B=0.2.
    \]
    What is the mean recurrence time of state $B$?

    \item In the insurance example, explain in one sentence why part (b) is not solved using $1/\pi_L$.
\end{enumerate}

\section{Answers}

\subsection*{Answer 1}

If
\[
|1-p-q|<1,
\]
then the number
\[
1-p-q
\]
lies between $-1$ and $1$.

Any number whose absolute value is less than 1 goes to 0 when raised to higher and higher powers.

Therefore,
\[
\lim_{n\to\infty}(1-p-q)^n=0.
\]

So the part of $P^n$ containing $(1-p-q)^n$ disappears in the long run.

\subsection*{Answer 2}

If the chain is ergodic, then every row of the limiting matrix is equal to $\pi$.

Since
\[
\pi=(0.3,0.7),
\]
we get
\[
\lim_{n\to\infty}P^n
=
\begin{pmatrix}
0.3 & 0.7\\
0.3 & 0.7
\end{pmatrix}.
\]

\subsection*{Answer 3}

The event that the first day is sunny can happen as:
\[
(S,S)
\]
or
\[
(S,R).
\]

Therefore,
\[
\Prob(S)=\Prob(S,S)+\Prob(S,R).
\]

So
\[
\Prob(S)=0.6+0.2=0.8.
\]

\subsection*{Answer 4}

The probability of leaving state $A$ in a given step is
\[
0.25.
\]

The time spent in $A$ before leaving has a geometric distribution with success probability $0.25$.

Therefore, the expected time spent in $A$ is
\[
\frac1{0.25}=4.
\]

\[
\boxed{4\text{ time steps}.}
\]

\subsection*{Answer 5}

For an ergodic Markov chain,
\[
\E_B[\tau_B]=\frac1{\pi_B}.
\]

Since
\[
\pi_B=0.2,
\]
we get
\[
\E_B[\tau_B]=\frac1{0.2}=5.
\]

\[
\boxed{5\text{ time steps}.}
\]

\subsection*{Answer 6}

Part (b) asks for the time until leaving low risk, while $1/\pi_L$ gives the average time between repeated visits to low risk in the long run.

\end{document}
